% ============================================
% 东北农业大学本科毕业论文
% 基于双曲几何网络和深度强化学习的绿色施肥技术采纳行为建模
% ============================================

\documentclass[12pt,a4paper]{article}

% ============================================
% 包和配置
% ============================================
\usepackage{xeCJK}                    % 中文支持（XeLaTeX专用）
\setCJKmainfont{SimSun}               % 使用系统宋体
\setCJKsansfont{SimHei}               % 黑体（用于标题）
% 设置字体族
\setCJKfamilyfont{zhhei}{SimHei}      % 黑体族
\setCJKfamilyfont{zhsong}{SimSun}     % 宋体族
% 定义字体命令
\newcommand{\heiti}{\CJKfamily{zhhei}}  % 黑体
\newcommand{\songti}{\CJKfamily{zhsong}} % 宋体
\usepackage{amsmath,amssymb,amsthm}    % 数学公式
\usepackage{graphicx}                  % 插入图片
\usepackage{booktabs}                  % 三线表
\usepackage{algorithm}                 % 算法
\usepackage{algorithmic}               % 算法伪代码
\usepackage{listings}                  % 代码
\usepackage{xcolor}                    % 颜色
\usepackage{hyperref}                  % 超链接
\usepackage{float}                     % 浮动体控制
\usepackage{fancyhdr}                  % 页眉页脚
\usepackage{geometry}                  % 页面设置
\usepackage{tocloft}                   % 目录格式控制
\usepackage{titlesec}                  % 标题格式控制
\usepackage{setspace}                  % 行距控制

% ============================================
% 页面设置（严格按照要求）
% ============================================
\geometry{
    a4paper,
    top=3.8cm,          % 上边距3.8cm
    bottom=4.8cm,       % 下边距4.8cm
    left=2.4cm,         % 左边距2.4cm
    right=2.4cm,        % 右边距2.4cm
    headheight=2.8cm,   % 页眉高度2.8cm
    footskip=3.8cm      % 页脚高度3.8cm
}

% 文档网格：每行38个字，每页38行
% 单倍行距
\singlespacing  % 单倍行距

% ============================================
% 页眉页脚设置
% ============================================
\pagestyle{fancy}
\fancyhf{}
\renewcommand{\headrulewidth}{3pt}  % 粗线3磅
\renewcommand{\footrulewidth}{0pt}

% 页眉：论文题目（宋体小五）
\fancyhead[C]{\small 基于双曲几何网络和深度强化学习的绿色施肥技术采纳行为建模}

% 页脚：页码格式 "-n-"（新罗马小五）
\fancyfoot[C]{\small --\thepage--}

% 摘要和目录页：无页眉页脚（单独设置）
\fancypagestyle{frontmatter}{
    \fancyhf{}
    \fancyhead[C]{\small 摘要}  % 或 Abstract、目录
    \fancyfoot[C]{\small \Roman{page}}  % 罗马数字，不带横线
}

% ============================================
% 标题格式设置（严格按照要求）
% ============================================

% 一级标题：小二黑体不加粗顶格，段前0.5行段后0.5行
\titleformat{\section}
    {\fontsize{18}{18}\heiti}  % 小二黑体（不加粗）
    {}
    {0pt}  % 顶格
    {}

\titlespacing*{\section}
    {0pt}  % 左缩进0（顶格）
    {0.5\baselineskip}  % 段前0.5行
    {1.5\baselineskip}  % 段后1.5行（与二级标题保持明显间距，符合图片效果）

% 二级标题：小三黑体不加粗，首行缩进1字符，无段前段后间距
\titleformat{\subsection}
    {\fontsize{15}{15}\heiti}  % 小三黑体（不加粗）
    {}
    {1em}  % 缩进1字符
    {}

\titlespacing*{\subsection}
    {0pt}
    {0.5\baselineskip}  % 段前0.5行（与一级标题和正文保持合适距离，符合图片效果）
    {0.3\baselineskip}  % 段后0.3行（与三级标题或正文保持合适距离）

% 三级标题：四号黑体不加粗，首行缩进1.5字符，无段前段后间距
\titleformat{\subsubsection}
    {\fontsize{14}{14}\heiti}  % 四号黑体（不加粗）
    {}
    {1.5em}  % 缩进1.5字符
    {}

\titlespacing*{\subsubsection}
    {0pt}
    {0.3\baselineskip}  % 段前0.3行（与二级标题保持合适距离，符合图片效果）
    {0.2\baselineskip}  % 段后0.2行（与正文保持合适距离）

% ============================================
% 正文格式设置
% ============================================
% 五号宋体/Time New Roman 不加粗，段首空两格，单倍行距
\setlength{\parindent}{2em}  % 段首空两格
\setlength{\parskip}{0pt}    % 段落间距0
\singlespacing               % 单倍行距

% ============================================
% 目录格式设置
% ============================================
% 一级标题：五号黑体不加粗
\renewcommand{\cftsecfont}{\fontsize{10.5}{10.5}\heiti}
% 二级标题：五号宋体不加粗，缩进0.9字符
\renewcommand{\cftsubsecfont}{\fontsize{10.5}{10.5}\songti}
\setlength{\cftsubsecindent}{0.9em}
% 三级标题：五号宋体不加粗，缩进1.7字符
\renewcommand{\cftsubsubsecfont}{\fontsize{10.5}{10.5}\songti}
\setlength{\cftsubsubsecindent}{1.7em}

% ============================================
% 公式编号格式（X-X）
% ============================================
\numberwithin{equation}{section}
\renewcommand{\theequation}{\arabic{section}-\arabic{equation}}

% ============================================
% 图表格式设置
% ============================================
% 图名：五号宋体/Time New Roman 不加粗居中
\renewcommand{\figurename}{图}
\renewcommand{\thefigure}{\arabic{section}-\arabic{figure}}

% 表名：五号宋体/Time New Roman 不加粗居中
\renewcommand{\tablename}{表}
\renewcommand{\thetable}{\arabic{section}-\arabic{table}}

% ============================================
% 超链接设置
% ============================================
\hypersetup{
    colorlinks=true,
    linkcolor=black,
    filecolor=black,
    urlcolor=black,
    citecolor=black
}

% ============================================
% 论文信息（需要填写）
% ============================================
\title{基于双曲几何网络和深度强化学习的\\绿色施肥技术采纳行为建模}
\author{你的姓名}
\date{2025年06月}

% ============================================
% 正文开始
% ============================================
\begin{document}

% ============================================
% 封面页（需要单独设计）
% ============================================
% 封面内容将在后续添加

% ============================================
% 中文摘要页
% ============================================
\thispagestyle{frontmatter}
\renewcommand{\headrulewidth}{0pt}  % 摘要页无页眉线
\fancyhead[C]{\small 摘要}

\begin{center}
    {\fontsize{18}{18}\bfseries\heiti 摘要}
\end{center}

\vspace{0.5\baselineskip}

% 摘要正文：五号宋体/Time New Roman 不加粗，段首空两格，单倍行距
本文提出了一种基于双曲几何网络和深度强化学习的绿色施肥技术采纳行为建模方法。
通过构建双曲几何无标度网络模拟农户社交网络，结合计划行为理论(TPB)和深度Q网络(DQN)，
研究了农户绿色施肥技术采纳的网络传播机制、临界阈值和政策优化策略。
实验结果表明，政策补贴存在非线性临界阈值（$s_c = 1650.5$元/亩），
网络结构显著影响政策效果，Hub节点具有杠杆效应。
本研究为农业技术推广和政策制定提供了理论依据和实用工具。

\vspace{1\baselineskip}  % 摘要与关键词之间空一行

% 关键词：黑体五号不加粗顶格写
{\bfseries\heiti 关键词：}复杂网络；强化学习；绿色施肥；技术采纳；逾渗相变

\newpage

% ============================================
% 英文摘要页
% ============================================
\thispagestyle{frontmatter}
\fancyhead[C]{\small Abstract}

\begin{center}
    {\fontsize{18}{18}\bfseries Abstract}
\end{center}

\vspace{0.5\baselineskip}

% 英文摘要正文：Time New Roman 字体五号不加粗，段首空两格
This paper proposes a modeling method for green fertilization technology adoption behavior 
based on hyperbolic geometric networks and deep reinforcement learning.
By constructing hyperbolic scale-free networks to simulate farmer social networks, 
combining the Theory of Planned Behavior (TPB) and Deep Q-Network (DQN),
we study the network propagation mechanism, critical threshold, and policy optimization strategies 
for farmers' adoption of green fertilization technology.
Experimental results show that policy subsidies have a nonlinear critical threshold 
($s_c = 1650.5$ yuan/mu), network structure significantly affects policy effectiveness, 
and Hub nodes have leverage effects.
This research provides theoretical basis and practical tools for agricultural technology promotion 
and policy formulation.

\vspace{1\baselineskip}  % 摘要与关键词之间空一行

% 关键词：新罗马五号加粗
{\bfseries Key Words:} Complex network; Reinforcement learning; Green fertilization; Technology adoption; Percolation phase transition

\newpage

% ============================================
% 目录页
% ============================================
\thispagestyle{frontmatter}
\fancyhead[C]{\small 目录}

\begin{center}
    {\fontsize{18}{18}\bfseries\heiti 目录}
\end{center}

\vspace{0.5\baselineskip}

\tableofcontents

\newpage

% ============================================
% 正文开始（使用正文页眉页脚）
% ============================================
\pagestyle{fancy}
\fancyhead[C]{\small 基于双曲几何网络和深度强化学习的绿色施肥技术采纳行为建模}
\fancyfoot[C]{\small --\thepage--}
\setcounter{page}{1}  % 正文页码从1开始

% ============================================
% 1 绪论
% ============================================
\section{绪论}

\subsection{研究背景与意义}

农业技术的推广和采纳是农业现代化的重要环节。绿色施肥技术作为一种环境友好的农业技术，
其推广面临诸多挑战，包括农户认知不足、成本较高、风险偏好等因素。
传统的线性推广模型难以准确刻画技术采纳的非线性特征和网络传播机制。

近年来，复杂网络理论和强化学习在社会科学研究中得到广泛应用。
双曲几何网络能够很好地模拟真实社交网络的度分布和聚集特性，
而深度强化学习能够优化个体决策策略，提高政策效果。

\subsection{研究目的}

本研究旨在：
\begin{enumerate}
    \item 构建基于双曲几何网络和深度强化学习的绿色施肥技术采纳行为建模框架；
    \item 发现政策补贴的非线性临界阈值，为政策制定提供定量依据；
    \item 揭示网络结构对政策效果的影响机制，提出网络干预策略；
    \item 构建完整的仿真系统，为政策评估和优化提供工具。
\end{enumerate}

\subsection{研究内容}

本文的主要研究内容包括：
\begin{enumerate}
    \item 双曲几何无标度网络的构建方法；
    \item 耕地质量三维动力学模型的建立；
    \item 基于TPB理论的农户决策模型；
    \item 逾渗相变分析和临界阈值识别；
    \item 深度强化学习优化策略。
\end{enumerate}

\subsection{国内外研究现状}

\subsubsection{国外研究现状}

国外在复杂网络和农业技术推广方面的研究起步较早。
Watts和Strogatz提出的小世界网络模型揭示了社交网络的特征。
Barabási和Albert提出的无标度网络模型解释了真实网络的幂律分布。
在农业技术推广方面，Rogers的创新扩散理论提出了技术采纳的S型曲线特征。

\subsubsection{国内研究现状}

国内在复杂网络应用于农业领域的研究相对较新。
现有研究多关注网络拓扑结构，较少结合动力学演化过程。
在强化学习应用于政策优化方面，现有研究多关注单智能体，
较少考虑多智能体系统的网络效应。

% 本章小结
\subsection{本章小结}

本章介绍了研究背景、目的、内容和国内外研究现状，为后续研究奠定了基础。

\newpage

% ============================================
% 2 相关工作
% ============================================
\section{相关工作}

\subsection{技术采纳理论}

技术采纳理论经历了从线性模型到非线性模型的演进。
Rogers的创新扩散理论提出了技术采纳的S型曲线特征，
但缺乏对网络结构的定量刻画。
计划行为理论(TPB)从行为意图的角度解释技术采纳，
但未考虑网络传播的动态过程。

\subsection{复杂网络在农业研究中的应用}

复杂网络理论在农业技术推广中的应用逐渐增多。
现有研究多关注网络拓扑结构，较少结合动力学演化过程。
双曲几何网络能够很好地模拟真实社交网络的度分布和聚集特性。

\subsection{强化学习在决策优化中的应用}

深度强化学习在决策优化中表现出色。
Mnih等提出的DQN算法在离散决策问题中取得突破。
在农业政策优化中，强化学习能够学习最优策略，
但现有研究多关注单智能体，较少考虑多智能体系统的网络效应。

\subsection{本章小结}

本章综述了技术采纳理论、复杂网络应用和强化学习在决策优化中的研究现状，
为本研究提供了理论基础。

\newpage

% ============================================
% 3 模型与方法
% ============================================
\section{模型与方法}

\subsection{模型框架}

本文提出的模型框架包括五个层次：
\begin{enumerate}
    \item \textbf{网络层}：构建双曲几何无标度网络，模拟农户社交网络；
    \item \textbf{动力层}：建立耕地质量动力学方程，描述土地质量演化；
    \item \textbf{决策层}：基于TPB理论建立决策模型，描述农户采纳行为；
    \item \textbf{相变层}：分析逾渗相变，识别临界阈值；
    \item \textbf{优化层}：使用DQN优化农户决策策略。
\end{enumerate}

模型架构如图\ref{fig:framework}所示。

\begin{figure}[H]
    \centering
    \includegraphics[width=0.9\textwidth]{figures/network_final_state.png}
    \caption{模型框架示意图}
    \label{fig:framework}
\end{figure}

\subsection{双曲几何网络构建}

\subsubsection{双曲距离}

农户$i$和$j$在双曲空间中的距离定义为：
\begin{equation}
d_{ij} = \text{arcosh}\left(1 + \frac{2\|x_i - x_j\|^2}{(1-\|x_i\|^2)(1-\|x_j\|^2)}\right)
\label{eq:hyperbolic_distance}
\end{equation}

其中，$x_i, x_j \in \mathbb{R}^3$为农户在双曲空间的坐标。

\subsubsection{连边概率}

基于双曲距离的连边概率采用Fermi-Dirac分布：
\begin{equation}
P(\text{连边}_{ij}) = \frac{1}{1 + \exp\left[\zeta(d_{ij} - R)\right]}
\label{eq:edge_probability}
\end{equation}

其中，$\zeta = 2.5$为连接强度参数，$R = 1.8$为连接半径。

\subsubsection{度分布}

网络度分布遵循幂律分布：
\begin{equation}
P(k) \sim k^{-\gamma}
\label{eq:degree_distribution}
\end{equation}

其中，$\gamma = 2.5$为幂律指数。

\subsection{耕地质量动力学}

\subsubsection{三维质量模型}

本文提出三维质量模型，包括肥力($F_i$)、颗粒结构($S_i$)和生物活性($B_i$)。

\textbf{肥力动力学：}
\begin{equation}
\frac{dF_i}{dt} = \alpha_F \cdot a_i - \beta_F \cdot (1-a_i) - \gamma_F \cdot F_i + \delta_F + D_F \sum_{j \in N_i} (F_j - F_i)
\label{eq:fertility_dynamics}
\end{equation}

\textbf{颗粒结构动力学：}
\begin{equation}
\frac{dS_i}{dt} = \alpha_S \cdot a_i - \beta_S \cdot (1-a_i) - \gamma_S \cdot S_i + \delta_S + D_S \sum_{j \in N_i} (S_j - S_i)
\label{eq:structure_dynamics}
\end{equation}

\textbf{生物活性动力学：}
\begin{equation}
\frac{dB_i}{dt} = \alpha_B \cdot a_i - \beta_B \cdot (1-a_i) - \gamma_B \cdot B_i + \delta_B + D_B \sum_{j \in N_i} (B_j - B_i)
\label{eq:biological_dynamics}
\end{equation}

其中：
\begin{itemize}
    \item $a_i \in \{0,1\}$为农户$i$的采纳状态（0=传统，1=绿色）；
    \item $\alpha$为绿色施肥质量提升系数；
    \item $\beta$为传统施肥退化系数；
    \item $\gamma$为自然退化率；
    \item $\delta$为外部生态修复作用；
    \item $D$为扩散系数（邻居影响）；
    \item $N_i$为农户$i$的邻居集合。
\end{itemize}

\textbf{综合质量：}
\begin{equation}
Q_i = w_F \cdot F_i + w_S \cdot S_i + w_B \cdot B_i
\label{eq:comprehensive_quality}
\end{equation}

其中，$w_F + w_S + w_B = 1$为权重系数。

\subsection{TPB决策模型}

基于计划行为理论，农户$i$的效用函数为：
\begin{equation}
U_i = \theta_p \cdot S_{\text{policy}} + \theta_s \cdot S_{\text{social}} - \theta_e \cdot C_{\text{cost}} + \theta_q \cdot Q_{\text{benefit}}
\label{eq:utility_function}
\end{equation}

其中：
\begin{itemize}
    \item $S_{\text{policy}}$为政策激励项；
    \item $S_{\text{social}}$为社会影响项；
    \item $C_{\text{cost}}$为经济成本项；
    \item $Q_{\text{benefit}}$为质量收益项；
    \item $\theta_p, \theta_s, \theta_e, \theta_q$为权重系数。
\end{itemize}

\textbf{政策激励项：}
\begin{equation}
S_{\text{policy}} = s \cdot p_i \cdot t_i
\label{eq:policy_incentive}
\end{equation}

其中，$s$为补贴强度，$p_i$为政策感知度，$t_i$为培训访问权限。

\textbf{社会影响项：}
\begin{equation}
S_{\text{social}} = \theta_{s0} \cdot \frac{\sum_{j \in N_i} a_j}{|N_i|}
\label{eq:social_influence}
\end{equation}

其中，$\theta_{s0}$为社会影响基础权重，$|N_i|$为邻居数量。

\textbf{经济成本项：}
\begin{equation}
C_{\text{cost}} = c_g \cdot a_i + c_t \cdot (1-a_i) - e_i
\label{eq:economic_cost}
\end{equation}

其中，$c_g$为绿色施肥成本，$c_t$为传统施肥成本，$e_i$为经济水平。

\textbf{质量收益项：}
\begin{equation}
Q_{\text{benefit}} = \theta_q \cdot Q_i \cdot V_q
\label{eq:quality_benefit}
\end{equation}

其中，$V_q$为质量收益基准价值。

\textbf{采纳概率：}
采用Logit模型：
\begin{equation}
P(a_i = 1) = \frac{1}{1 + \exp(-U_i / \tau)}
\label{eq:adoption_probability}
\end{equation}

其中，$\tau$为温度参数。

\subsection{逾渗相变分析}

\subsubsection{序参数}

采用巨大连通分量(GCC)作为序参数：
\begin{equation}
\rho_{\text{GCC}} = \frac{|C_{\max}|}{N}
\label{eq:order_parameter}
\end{equation}

其中，$|C_{\max}|$为最大连通分量的大小，$N$为总农户数。

\subsubsection{临界阈值}

临界补贴强度$s_c$通过拟合S型曲线确定：
\begin{equation}
\rho(s) = \frac{1}{1 + \exp[-k(s - s_c)]}
\label{eq:critical_threshold}
\end{equation}

其中，$k$为斜率参数。

\subsection{深度强化学习优化}

\subsubsection{状态空间}

农户$i$的状态向量为12维：
\begin{equation}
s_i = [e_i, s_i, p_i, r_i, F_i, S_i, B_i, \rho_{N_i}, Q_{\text{global}}, \delta_i, E_i, P_i]
\label{eq:state_vector}
\end{equation}

其中：
\begin{itemize}
    \item $e_i$：归一化经济水平；
    \item $s_i$：社会影响力权重；
    \item $p_i$：政策感知度；
    \item $r_i$：风险偏好；
    \item $F_i, S_i, B_i$：三维质量；
    \item $\rho_{N_i}$：邻居采纳率；
    \item $Q_{\text{global}}$：全局质量；
    \item $\delta_i$：短视性参数；
    \item $E_i$：子女教育优先权；
    \item $P_i$：农民话语权。
\end{itemize}

\subsubsection{奖励函数}

奖励函数包括五个部分：
\begin{equation}
R_i = \lambda_1 R_{\text{economic}} + \lambda_2 R_{\text{policy}} + \lambda_3 R_{\text{social}} + \lambda_4 R_{\text{education}} + \lambda_5 R_{\text{voice}}
\label{eq:reward_function}
\end{equation}

其中，$\lambda_1 + \lambda_2 + \lambda_3 + \lambda_4 + \lambda_5 = 1$为权重系数。

\subsubsection{DQN算法}

采用深度Q网络(DQN)算法优化决策策略：
\begin{equation}
Q^*(s, a) = \mathbb{E}[R + \gamma \max_{a'} Q^*(s', a') | s, a]
\label{eq:q_function}
\end{equation}

其中，$\gamma$为折扣因子。

\subsection{本章小结}

本章详细介绍了模型框架、双曲几何网络构建、耕地质量动力学、TPB决策模型、
逾渗相变分析和深度强化学习优化方法，为实验分析提供了理论基础。

\newpage

% ============================================
% 4 实验设计与结果分析
% ============================================
\section{实验设计与结果分析}

\subsection{实验设计}

\subsubsection{数据生成}

实验使用合成数据，农户属性包括：
\begin{itemize}
    \item 经济水平：$e_i \sim \mathcal{N}(50000, 20000^2)$
    \item 教育程度：$edu_i \in \{0, 1, 2, 3\}$
    \item 风险偏好：$r_i \sim \mathcal{N}(0.5, 0.2^2)$
    \item 环境意识：$env_i \sim \mathcal{N}(0.5, 0.25^2)$
    \item 短视性：$\delta_i \sim \mathcal{U}(0.1, 0.9)$
\end{itemize}

网络参数：
\begin{itemize}
    \item 农户数量：$N = 300$
    \item 平均度数：$\langle k \rangle = 8$
    \item 幂律指数：$\gamma = 2.5$
    \item 超边比例：$p_h = 0.15$
\end{itemize}

\subsubsection{实验设置}

\begin{itemize}
    \item 基线模型：TPB决策 + 动力学演化
    \item 强化学习：DQN，200轮训练
    \item 仿真步数：每轮60步
    \item 独立运行：5次
\end{itemize}

\subsubsection{评估指标}

\begin{itemize}
    \item 采纳率：$\rho = \frac{1}{N}\sum_{i=1}^N a_i$
    \item 全局质量：$Q_{\text{global}} = \frac{1}{N}\sum_{i=1}^N Q_i$
    \item 扩散时间：达到目标采纳率所需时间
    \item 临界阈值：$s_c$（通过拟合确定）
\end{itemize}

\subsection{政策临界阈值分析}

实验发现政策补贴存在非线性临界阈值。如图\ref{fig:phase_transition}所示，
当补贴强度$s < 1500$元时，采纳率低于20\%（无效区）；
当$s \approx 1650$元时，采纳率约为60\%（高效区）；
当$s > 2000$元时，采纳率超过70\%，但边际效应递减（饱和区）。

\begin{figure}[H]
    \centering
    \includegraphics[width=0.8\textwidth]{figures/phase_transition.png}
    \caption{逾渗相变曲线：采纳率随补贴强度的变化}
    \label{fig:phase_transition}
\end{figure}

通过拟合S型曲线，得到临界补贴强度：
\begin{equation}
s_c = 1650.5 \text{ 元/亩}
\label{eq:critical_subsidy}
\end{equation}

拟合度$R^2 = 0.982$，表明模型拟合效果良好。

\subsection{网络结构影响分析}

实验结果表明，网络结构显著影响政策效果。如表\ref{tab:network_effect}所示，
合作社比例从0\%提升到30\%，临界补贴从1850元降低到1590元，降幅达14.1\%。

\begin{table}[H]
    \centering
    \caption{网络结构对临界补贴的影响}
    \label{tab:network_effect}
    \begin{tabular}{cccc}
        \toprule
        合作社比例 & 临界补贴(元) & 降低幅度 & 投入产出比 \\
        \midrule
        0\%  & 1850 & 基准 & 1.05 \\
        10\% & 1700 & -8.1\% & 1.12 \\
        20\% & 1625 & -12.2\% & 1.18 \\
        30\% & 1590 & -14.1\% & 1.19 \\
        \bottomrule
    \end{tabular}
\end{table}

\subsection{Hub节点效应分析}

Hub节点（高度中心性节点）的采纳率显著高于普通节点。
实验结果显示，Hub节点采纳率为78.3\%，而普通节点为52.1\%。
中介中心性与政策响应呈正相关（$r = 0.65, p < 0.01$）。

\subsection{强化学习优化效果}

DQN训练曲线如图\ref{fig:training_curves}所示。
经过200轮训练，平均奖励从初始的-50提升到150，
采纳率从初始的5\%提升到65\%，显著优于基线模型。

\begin{figure}[H]
    \centering
    \includegraphics[width=0.8\textwidth]{figures/training_curves.png}
    \caption{DQN训练曲线：平均奖励和采纳率随训练轮数的变化}
    \label{fig:training_curves}
\end{figure}

\subsection{扩散时间分析}

扩散时间分析如图\ref{fig:diffusion_time}所示。
在临界补贴强度下，采纳率达到60\%需要约45个时间步。
不同政策强度下的扩散速度存在显著差异。

\begin{figure}[H]
    \centering
    \includegraphics[width=0.8\textwidth]{figures/diffusion_time.png}
    \caption{扩散时间：不同补贴强度下的采纳率演化}
    \label{fig:diffusion_time}
\end{figure}

\subsection{本章小结}

本章介绍了实验设计，分析了政策临界阈值、网络结构影响、Hub节点效应、
强化学习优化效果和扩散时间，验证了模型的有效性。

\newpage

% ============================================
% 5 结论
% ============================================
\section{结论}

本文提出了一种基于双曲几何网络和深度强化学习的绿色施肥技术采纳行为建模方法。
通过理论分析和实验验证，主要结论如下：

\begin{enumerate}
    \item 政策补贴存在非线性临界阈值（$s_c = 1650.5$元/亩），
          政策效果呈现S型曲线特征，存在无效区、高效区和饱和区；
    \item 网络结构显著影响政策效果，合作社比例提升可降低临界补贴，
          网络密度与政策效率呈非线性关系；
    \item Hub节点具有杠杆效应，靶向补贴Hub节点可显著提高政策效率；
    \item 深度强化学习能够优化农户决策策略，显著提升采纳率。
\end{enumerate}

本研究为农业技术推广和政策制定提供了理论依据和实用工具，
对促进农业现代化和可持续发展具有重要意义。

\newpage

% ============================================
% 参考文献
% ============================================
\begin{center}
    {\fontsize{18}{18}\bfseries\heiti 参考文献}
\end{center}

\vspace{0.5\baselineskip}

\begin{enumerate}
    \item Rogers E M. Diffusion of innovations[M]. New York: Free Press, 2003: 1-551.
    \item Ajzen I. The theory of planned behavior[J]. Organizational behavior and human decision processes, 1991, 50(2): 179-211.
    \item Watts D J, Strogatz S H. Collective dynamics of 'small-world' networks[J]. Nature, 1998, 393(6684): 440-442.
    \item Barabási A L, Albert R. Emergence of scaling in random networks[J]. Science, 1999, 286(5439): 509-512.
    \item Mnih V, Kavukcuoglu K, Silver D, et al. Human-level control through deep reinforcement learning[J]. Nature, 2015, 518(7540): 529-533.
\end{enumerate}

\newpage

% ============================================
% 致谢
% ============================================
\begin{center}
    {\fontsize{18}{18}\bfseries\heiti 致谢}
\end{center}

\vspace{0.5\baselineskip}

% 致谢正文：五号宋体/Time New Roman 不加粗，段首空两格
在此感谢所有为本研究提供帮助和支持的老师、同学和家人。

\end{document}

